\begin{center}
\begin{tikzpicture}[font=\small,>=stealth]
\draw[gray!40,rounded corners] (0,0) rectangle (14.2,9);
\node[font=\small\bfseries] at (7.1,8.35) {防混淆：内分线沿切线；顶点处三角形退化};

\begin{scope}[shift={(3.6,4.7)},scale=.4]

\coordinate (L) at (-5,0); \coordinate (R) at (5,0);
\coordinate (P) at (3.75,3); \coordinate (M) at (2.4,0);
\draw[blue!75!black,thick] (L)--(P);
\draw[orange!85!black,thick] (P)--(R);
\draw[gray] (L)--(R);
\draw[magenta!80!black,very thick] (M)--(P);
\foreach \q in {L,R,P,M} \fill (\q) circle (3pt);
\node[below left] at (L) {$F_1$}; \node[below right] at (R) {$F_2$};
\node[above left] at (P) {$P$}; \node[below left] at (M) {$M$};

\draw[magenta!80!black] (1.7,-1.5556)--(4.9,5.5556);
\draw[gray,dashed] (0,4.6875)--(7,1.5375);
\node[above] at (6,2.1) {法线};
\end{scope}
\node at (3.7,2.5) {紫实线：内角平分线／切线};
\node at (3.7,1.6) {灰虚线：外角平分线／法线};
\begin{scope}[shift={(10.8,4.7)},scale=.42]
\draw[gray] (-5.8,0)--(6,0);
\draw[teal!75!black,domain=-.55:.55,samples=60,variable=\t] plot ({3*cosh(\t)},{4*sinh(\t)});
\draw[magenta!80!black] (3,-3)--(3,3);
\fill (-5,0) circle(4pt) node[below] {$F_1$};
\fill (5,0) circle(4pt) node[below] {$F_2$};
\fill (3,0) circle(4pt);\node[above left] at (3,0) {$P=A$};
\end{scope}
\node at (10.7,2.5) {$P,F_1,F_2$ 共线：没有三角形};
\node at (10.7,1.6) {顶点切线 $x=a$ 仍存在};
\node at (7.1,.65) {必须保留条件 $y_0\ne0$，才能谈焦点三角形的顶角内分线。};

\end{tikzpicture}
\end{center}
